% Generated draft: requires mathematical and source review.
\documentclass[11pt]{article}
\usepackage[margin=25mm]{geometry}
\usepackage{fontspec}
\setmainfont{Latin Modern Roman}
\usepackage{amsmath,amssymb,mathtools,unicode-math}
\setmathfont{Latin Modern Math}
\usepackage{enumitem,needspace,xurl,hyperref}
\setlength{\parindent}{0pt}
\setlength{\parskip}{5pt}
\setlength{\emergencystretch}{3em}
\newcommand{\N}{\mathbb{N}}
\newcommand{\Z}{\mathbb{Z}}
\newcommand{\Q}{\mathbb{Q}}
\newcommand{\R}{\mathbb{R}}
\newcommand{\C}{\mathbb{C}}
\newcommand{\F}{\mathbb{F}}
\newcommand{\A}{\mathbb{A}}
\newcommand{\PP}{\mathbb{P}}
\newcommand{\E}{\mathbb{E}}
\newcommand{\eps}{\varepsilon}
\newcommand{\dd}{\,\mathrm{d}}
\begin{document}
\section*{UnsolvedMath ID: 20001951}
\textbf{Source problem code:} catalogue entry
\subsection{Definitions and mathematical statement}
Let $S$ be a complete compact 3-Sasakian manifold; its metric cone is hyperk\"ahler. Let a connected compact Lie group $G\subset\operatorname{Aut}(S)$ preserve the 3-Sasakian structure. For contact forms $\eta_1,\eta_2,\eta_3$, write the contact moment map as
\[
 \mu:S\longrightarrow\mathfrak g^*\otimes\R^3,
 \qquad \langle\mu_i(s),A\rangle=\eta_i(A^\#)_s,
 \qquad Z=\mu^{-1}(0),\quad S_{\mathrm{red}}=Z/G.
\]
If zero is regular and the $G$-action on $Z$ is locally free, use ordinary rational or real cohomology of the quotient orbifold. Restriction and the rational/real identification $H_G^*(Z;\mathbb K)\cong H^*(Z/G;\mathbb K)$ define
\[
 \kappa_S:H_G^*(S;\mathbb K)\longrightarrow
                  H^*(S_{\mathrm{red}};\mathbb K),
 \qquad \mathbb K\in\{\Q,\R\}.
\]
Here $H_G^*$ denotes Borel equivariant cohomology. For singular quotients, an alternative target theory is not specified by the original conjecture.

\textbf{Conjecture 2.1, C.~Boyer.} This map should be surjective up to the middle dimension of the reduced space. For an odd-dimensional quotient of dimension $4m-1$, the inclusive lower-middle interpretation is
\[
 \kappa_S^q\text{ is onto for }0\le q\le2m-1.
\]
The source does not formally resolve endpoint conventions or general regularity assumptions. Boyer's Comment 2.2 reports the conjecture for circle reductions of spheres and says that in odd dimension only a middle-dimensional range can be expected. It cites examples of Boyer, Galicki and Piccinni [BGP02] whose cohomology computations motivate the question.

\textbf{Sphere case and proposed hyperk\"ahler route.} In Comment 2.3, Boyer and T.~Hausel warn that the nonabelian case may need extra hypotheses and discuss abelian reductions. For $G\subset\operatorname{Sp}(n+1)$ acting on $S^{4n+3}$, they conjecture the same range when $L=S^{4n+3}\mathbin{///}G$ is an orbifold. In the regular locally free case, $\dim L=4(n+1-\dim G)-1$, so $m=n+1-\dim G$. Their tentative strategy uses the corresponding hyperk\"ahler reduction of $\mathbb H^{n+1}$ at a suitable central parameter with orbifold quotient and a hypercompactness hypothesis ensuring a middle-dimensional core. More precisely a moment-map parameter has type $\xi\in(\mathfrak g^*\otimes\R^3)^G$, not an element of the group center itself. Definition 6.9 of hypercompactness is not supplied. The proposed transfer is not assumed automatic.

\textbf{Supplied conditional boundary-transfer statement.} Suppose $L$ is the boundary of a compact oriented filling $W$ of dimension $4m$ for which Poincar\'e--Lefschetz duality over $\mathbb K$ holds. Assume $W$ retracts onto a core of dimension at most $2m$, and there is a graded map
\[
 \kappa_W:H_G^*(S^{4n+3};\mathbb K)\longrightarrow H^*(W;\mathbb K)
\]
onto in degrees at most $2m-1$, satisfying $i^*\kappa_W=\kappa_S$ for $i:L\hookrightarrow W$. This explicitly states the aid's Kirwan-compatible filling assumption. Define
\[
 I:H_{2m}(W;\mathbb K)\xrightarrow{\mathrm{PD}^{-1}}
       H^{2m}(W,L;\mathbb K)\longrightarrow H^{2m}(W;\mathbb K).
\]
The aid asserts surjectivity of $\kappa_S^q$ for $q\le2m-2$ and a natural identification
\[
 \operatorname{coker}\kappa_S^{2m-1}\cong\ker I.
\]
Thus the inclusive endpoint $2m-1$ additionally requires injectivity of the middle intersection map. A middle-dimensional core alone does not provide it. The aid's abstract example $W=C\times D^{2m}$, for a closed connected oriented $2m$-manifold $C$, has such a core but a zero middle intersection map with nonzero domain. This is an obstruction to the abstract boundary-transfer argument; it is not asserted to be an actual 3-Sasakian reduction or a counterexample to Boyer's conjecture. Existence of a compatible filling with these hypotheses remains separate from hypercompactness and from hyperk\"ahler Kirwan surjectivity.
\subsection{Short English statement}
Boyer conjectures that 3-Sasakian Kirwan maps are surjective through the lower middle degree. The supplied aid gives a conditional transfer from a compatible filling and identifies the endpoint obstruction as a middle intersection kernel. A middle-dimensional core alone does not remove that obstruction or settle the conjecture.
\subsection{Sources}
\begin{itemize}
\item Local UnsolvedMath registry, permanent ID 20001951, source code catalogue entry; canonical \texttt{aim-geometry-notes.json}, notes[288].
\item AIM workshop \emph{Moment maps and surjectivity in various geometries}, Conjecture 2.1 (C.~Boyer), Comment 2.2 (C.~Boyer), Comment 2.3 (C.~Boyer and T.~Hausel): \url{https://aimath.org/WWN/momentmaps/momentmaps.pdf}.
\item Supplied reference [BGP02], attributed to Boyer, Galicki and Piccinni. Full bibliographic entry and the cross-referenced Definition 6.9 of hypercompactness are absent. The source describes Boyer and Hausel's intended post-workshop investigation, not an accomplished proof.
\item Registry research attempt 1, dated 2026-08-14: conditional result, status at run open, rigor ``proved'', novelty confidence low. The machine-generated, non-peer-reviewed aid requires independent verification before citation.
\end{itemize}
\end{document}
